\documentclass[14pt,a4paper,oneside]{extbook}
\usepackage[utf8]{vietnam}
\usepackage{amsmath,amssymb,mathrsfs,fancyhdr,enumerate,multirow,makecell}
\usepackage{tkz-euclide,tikz-3dplot,tikz,tkz-tab}
\usetikzlibrary{arrows.meta,calc,intersections,angles,quotes,backgrounds,shapes.geometric,patterns,through,shadows,positioning}
\usepackage{yhmath}
\usepackage[top=0.7cm, bottom=0.7cm, left=1.0cm, right=1.0cm, headheight=16pt, headsep=3mm, footskip=5mm]{geometry}
\usepackage[hidelinks,unicode]{hyperref}
\usepackage{qrcode}
\usepackage[many]{tcolorbox}
\usepackage{fontawesome}

\pagestyle{fancy}
\fancyhf{}
\fancyhead[L]{\footnotesize\bfseries\color{teal!85!black}\faBook~SỔ TAY KIẾN THỨC TOÁN 4}
\fancyhead[R]{\footnotesize\bfseries\color{blue!85!black}Lớp Toán Thầy Lĩnh -- loptoanthaylinh.vn}
\fancyfoot[C]{\footnotesize\bfseries\color{gray!80!black}-- Sổ tay kiến thức: Bài 49. Dãy số liệu thống kê --}
\renewcommand{\headrulewidth}{0.4pt}

\begin{document}

% --- BANNER TIÊU ĐỀ BÀI HỌC ---
\noindent
\begin{tcolorbox}[
  enhanced,
  colback=teal!8!white,
  colframe=teal!80!black,
  arc=3mm, boxrule=1pt,
  left=3mm, right=3mm, top=2mm, bottom=2mm,
  halign=center
]
  {\Large\bfseries\color{teal!90!black}\faBookmark\ BÀI 49: DÃY SỐ LIỆU THỐNG KÊ}\\[1.5pt]
  {\footnotesize\itshape\color{black!75} Khái niệm dãy số liệu, cách sắp xếp, phân tích thông tin và tìm giá trị đặc trưng trên 1 trang A4}
\end{tcolorbox}

\vspace{2mm}

% --- KHỐI 1 & 2: LÝ THUYẾT & KĨ NĂNG ---
\noindent
\begin{minipage}[t]{0.49\textwidth}
  \begin{tcolorbox}[
    enhanced,
    title={\footnotesize\bfseries\color{white}\faDatabase\ 1. Khái niệm dãy số liệu},
    colback=blue!2!white,
    colframe=blue!75!black,
    arc=2.5mm, boxrule=0.8pt, top=1.8mm, bottom=1.8mm, left=2mm, right=2mm
  ]
    \footnotesize
    $\bullet$ Khi thu thập, kiểm đếm các số liệu và viết các số đó \textbf{theo thứ tự}, ta được một \textbf{dãy số liệu thống kê}.\\
    $\bullet$ \textbf{Ví dụ}: Dãy số đo chiều cao (cm) của 4 bạn:\\
    \begin{center}
      $132;\ 135;\ 138;\ 140$
    \end{center}
    --- Số thứ nhất là $132$; số thứ hai là $135$; số thứ ba là $138$; số thứ tư là $140$.\\
    --- Dãy này có tất cả \textbf{4 số liệu}.
  \end{tcolorbox}
\end{minipage}%
\hfill
\begin{minipage}[t]{0.49\textwidth}
  \begin{tcolorbox}[
    enhanced,
    title={\footnotesize\bfseries\color{white}\faSearch\ 2. Khai thác thông tin từ dãy},
    colback=teal!2!white,
    colframe=teal!75!black,
    arc=2.5mm, boxrule=0.8pt, top=1.8mm, bottom=1.8mm, left=2mm, right=2mm
  ]
    \footnotesize
    $\bullet$ Quan sát dãy số liệu, ta có thể biết:
    \par
    --- Dãy có \textbf{bao nhiêu số liệu}.\\
    --- Số \textbf{lớn nhất (Max)}, số \textbf{bé nhất (Min)}.\\
    --- Sắp xếp dãy theo thứ tự: từ bé đến lớn hoặc từ lớn đến bé.\\
    --- So sánh, tính chênh lệch giữa các số liệu.\\
    --- Tính tổng hoặc tính \textbf{số trung bình cộng}.
  \end{tcolorbox}
\end{minipage}

\vspace{2mm}

% --- KHỐI 3: SƠ ĐỒ TƯ DUY TIKZ ---
\noindent
\begin{tcolorbox}[
  enhanced,
  title={\footnotesize\bfseries\color{white}\faSitemap\ SƠ ĐỒ TƯ DUY: TIẾN TRÌNH XỬ LÍ DÃY SỐ LIỆU THỐNG KÊ},
  colback=yellow!2!white,
  colframe=orange!85!black,
  arc=2.5mm, boxrule=0.8pt, top=1.5mm, bottom=1.5mm, left=1mm, right=1mm,
  halign=center
]
  \begin{tikzpicture}[
    scale=0.88,
    card/.style={rectangle, rounded corners=2mm, draw=teal!80!black, fill=teal!10, thick, inner sep=3pt, font=\scriptsize\bfseries, align=center, text width=3.8cm},
    arr/.style={-{Stealth[scale=0.85]}, thick, draw=black!70, shorten >= 2.5pt, shorten <= 2.5pt}
  ]
    \node[card, draw=blue!75!black, fill=blue!10] (c1) at (0, 0) {\textbf{1. THU THẬP}\\Ghi chép số liệu\\theo thứ tự thực tế};
    \node[card, draw=orange!85!black, fill=orange!10] (c2) at (4.8, 0) {\textbf{2. ĐẾM SỐ LƯỢNG}\\Xác định số thứ nhất,\\thứ hai... đến hết};
    \node[card, draw=purple!80!black, fill=purple!10] (c3) at (9.6, 0) {\textbf{3. SẮP XẾP SỐ LIỆU}\\$\bullet$ Từ bé đến lớn\\$\bullet$ Từ lớn đến bé};
    \node[card, draw=teal!80!black, fill=teal!10] (c4) at (14.4, 0) {\textbf{4. TRÍCH XUẤT}\\$\bullet$ Số lớn nhất, bé nhất\\$\bullet$ TBC \& Chênh lệch};

    \draw[arr] (c1) -- (c2);
    \draw[arr] (c2) -- (c3);
    \draw[arr] (c3) -- (c4);
  \end{tikzpicture}
\end{tcolorbox}

\vspace{2mm}

% --- KHỐI 4: VÍ DỤ MINH HỌA & BẪY TRẮC NGHIỆM ---
\noindent
\begin{minipage}[t]{0.53\textwidth}
  \begin{tcolorbox}[
    enhanced,
    title={\footnotesize\bfseries\color{white}\faCheckCircle\ Ví dụ thực hành chuẩn mực},
    colback=green!2!white,
    colframe=green!60!black,
    arc=2.5mm, boxrule=0.8pt, top=1.8mm, bottom=1.8mm, left=2.2mm, right=2.2mm
  ]
    \footnotesize
    \textbf{Cho dãy số điểm kiểm tra của 5 bạn}:\\
    \begin{center}
      $8;\ 10;\ 7;\ 9;\ 8$
    \end{center}
    --- Dãy trên có $5$ số liệu.\\
    --- Số thứ nhất là $8$, số thứ hai là $10$, số thứ ba là $7$..\\
    --- Điểm cao nhất: $10$; Điểm thấp nhất: $7$.\\
    --- Sắp xếp từ bé đến lớn: $7;\ 8;\ 8;\ 9;\ 10$.\\
    --- Điểm trung bình: $(7 + 8 + 8 + 9 + 10) : 5 = \mathbf{8.4}$.
  \end{tcolorbox}
\end{minipage}%
\hfill
\begin{minipage}[t]{0.45\textwidth}
  \begin{tcolorbox}[
    enhanced,
    title={\footnotesize\bfseries\color{white}\faExclamationTriangle\ Mẹo tránh bẫy trắc nghiệm},
    colback=red!2!white,
    colframe=red!70!black,
    arc=2.5mm, boxrule=0.8pt, top=1.8mm, bottom=1.8mm, left=2.2mm, right=2.2mm
  ]
    \footnotesize
    $\bullet$ \textbf{Bẫy phân biệt "số thứ mấy" và "giá trị"}:\\
    --- "Số thứ ba là số nào?": Trả lời là $7$.\\
    --- "Có bao nhiêu số 8?": Trả lời là có 2 số.\\
    $\bullet$ \textbf{Bẫy đếm sót số liệu lặp lại}:\\
    Khi sắp xếp lại dãy số liệu, các số bằng nhau vẫn phải viết đầy đủ (ví dụ có 2 số 8 thì viết cả 2 số 8).
  \end{tcolorbox}
\end{minipage}

\end{document}
